\section{Step 3: Basic Argument Parsing}

Real clients accept flags before the URL. We introduce a straightforward
\texttt{for} loop over \texttt{argv} rather than \texttt{getopt\_long}, so
every branch is visible and easy to extend in later steps.

\subsection{Why a Manual Loop?}

\begin{itemize}[leftmargin=1.4em]
    \item No extra dependencies or platform quirks
    \item Optional values (e.g.\ \texttt{-T} with or without a proxy) are
          easier to express than with classic \texttt{getopt}
    \item Students see exactly when the index advances past an option argument
\end{itemize}

\subsection{State We Introduce}

\begin{itemize}[leftmargin=1.4em]
    \item \texttt{url\_arg} --- the single positional URL (required)
    \item \texttt{user\_agent} --- default browser-like string, overridable
          with \texttt{-A} / \texttt{--user-agent}
    \item Placeholders for Tor and redirect policy (wired up in later steps)
\end{itemize}

We also expand \texttt{usage()} so it documents the options we claim to
support. Empty stubs (\texttt{use\_tor}, \texttt{follow\_redirects}, \ldots)
document the roadmap without implementing behavior yet.

\subsection{New Code --- Expanded Usage}

\begin{lstlisting}[style=snippet,language=C]
static inline void usage(const char *progname)
{
    fprintf(stderr,
        "Usage: %s [options] <URL>\n\n"
        "Options:\n"
        "  -A, --user-agent STR     Set User-Agent\n"
        "  -T, --tor [proxy]        Use Tor (default 127.0.0.1:9050)\n"
        "  -H, --header STR         Add custom header (repeatable)\n"
        "  -d, --data DATA          POST data (@file supported)\n"
        "  -o, --output FILE        Write output to file\n"
        "  -u, --user USER:PASS     Basic authentication\n"
        "  -I, --head               HEAD request\n"
        "      --no-location        Do not follow redirects\n"
        "      --max-redirects N    Max redirects (default: 10)\n",
        progname);
    exit(EXIT_FAILURE);
}
\end{lstlisting}

\subsection{New Code --- Parsing Loop (core options)}

For this step we fully implement \texttt{-A} and recognition of the URL; Tor
and the other flags are parsed enough to set variables we will use later.

\begin{lstlisting}[style=snippet,language=C]
const char *user_agent =
    "Mozilla/5.0 (X11; Linux x86_64; rv:128.0) "
    "Gecko/20100101 Firefox/128.0";
const char *url_arg = NULL;

for (int i = 1; i < argc; i++) {
    if (strcmp(argv[i], "-A") == 0 ||
        strcmp(argv[i], "--user-agent") == 0) {
        if (i + 1 < argc) user_agent = argv[++i];
        else usage(argv[0]);
    }
    else if (url_arg == NULL) {
        url_arg = argv[i];
    }
}

if (url_arg == NULL)
    usage(argv[0]);
\end{lstlisting}

\textbf{Rule of thumb:} option switches that take a value always
\texttt{++i} after validating \texttt{i + 1 < argc}. Anything that is not a
recognized option becomes the URL---first one wins.

\FullSourceStart{urlgetstep03}
\begin{lstlisting}[style=oldcode,language=C,firstnumber=1]
/* urlget.c -- Step 3: argument parsing */
#include <stdio.h>
#include <stdlib.h>
\end{lstlisting}
\begin{lstlisting}[style=newcode,language=C,firstnumber=last]
#include <string.h>
#include <stdbool.h>
\end{lstlisting}
\begin{lstlisting}[style=oldcode,language=C,firstnumber=last]

\end{lstlisting}
\begin{lstlisting}[style=newcode,language=C,firstnumber=last]
static inline void usage(const char *progname)
{
    fprintf(stderr,
        "Usage: %s [options] <URL>\n\n"
        "Options:\n"
        "  -A, --user-agent STR     Set User-Agent\n"
        "  -T, --tor [proxy]        Use Tor (default 127.0.0.1:9050)\n"
        "  -H, --header STR         Add custom header (repeatable)\n"
        "  -d, --data DATA          POST data (@file supported)\n"
        "  -o, --output FILE        Write output to file\n"
        "  -u, --user USER:PASS     Basic authentication\n"
        "  -I, --head               HEAD request\n"
        "      --no-location        Do not follow redirects\n"
        "      --max-redirects N    Max redirects (default: 10)\n",
        progname);
    exit(EXIT_FAILURE);
}
\end{lstlisting}
\begin{lstlisting}[style=oldcode,language=C,firstnumber=last]
int main(int argc, char *argv[])
{
\end{lstlisting}
\begin{lstlisting}[style=newcode,language=C,firstnumber=last]
    const char *user_agent =
        "Mozilla/5.0 (X11; Linux x86_64; rv:128.0) "
        "Gecko/20100101 Firefox/128.0";
    const char *default_tor = "127.0.0.1:9050";

    const char *url_arg = NULL;
    const char *tor_proxy = NULL;
    bool use_tor = false;
    bool follow_redirects = true;
    int max_redirects = 10;
    bool head_request = false;
    char extra_headers[4096] = {0};
    char post_data[8192] = {0};
    char output_file[256] = {0};
    char auth_header[256] = {0};

    for (int i = 1; i < argc; i++) {
        if (strcmp(argv[i], "-A") == 0 ||
            strcmp(argv[i], "--user-agent") == 0) {
            if (i + 1 < argc) user_agent = argv[++i];
            else usage(argv[0]);
        }
        else if (strcmp(argv[i], "-T") == 0 ||
                 strcmp(argv[i], "--tor") == 0) {
            use_tor = true;
            if (i + 1 < argc && argv[i + 1][0] != '-' &&
                strchr(argv[i + 1], ':')) {
                tor_proxy = argv[i + 1];
                i++;
            }
        }
        else if (strcmp(argv[i], "-H") == 0 ||
                 strcmp(argv[i], "--header") == 0) {
            if (i + 1 < argc) {
                strncat(extra_headers, argv[++i],
                        sizeof(extra_headers) -
                        strlen(extra_headers) - 3);
                strcat(extra_headers, "\r\n");
            } else usage(argv[0]);
        }
        else if (strcmp(argv[i], "-d") == 0 ||
                 strcmp(argv[i], "--data") == 0) {
            if (i + 1 < argc) {
                const char *val = argv[++i];
                if (val[0] == '@') {
                    FILE *f = fopen(val + 1, "r");
                    if (f) {
                        fread(post_data, 1,
                              sizeof(post_data) - 1, f);
                        fclose(f);
                    }
                } else {
                    strncpy(post_data, val,
                            sizeof(post_data) - 1);
                }
            } else usage(argv[0]);
        }
        else if (strcmp(argv[i], "-o") == 0 ||
                 strcmp(argv[i], "--output") == 0) {
            if (i + 1 < argc)
                strncpy(output_file, argv[++i],
                        sizeof(output_file) - 1);
            else usage(argv[0]);
        }
        else if (strcmp(argv[i], "-u") == 0 ||
                 strcmp(argv[i], "--user") == 0) {
            if (i + 1 < argc)
                snprintf(auth_header, sizeof(auth_header),
                    "Authorization: Basic %s\r\n", argv[++i]);
            else usage(argv[0]);
        }
        else if (strcmp(argv[i], "-I") == 0 ||
                 strcmp(argv[i], "--head") == 0) {
            head_request = true;
        }
        else if (strcmp(argv[i], "--no-location") == 0) {
            follow_redirects = false;
        }
        else if (strcmp(argv[i], "--max-redirects") == 0 &&
                 i + 1 < argc) {
            max_redirects = atoi(argv[++i]);
        }
        else if (url_arg == NULL) {
            url_arg = argv[i];
        }
    }

    if (url_arg == NULL)
        usage(argv[0]);
    if (use_tor && tor_proxy == NULL)
        tor_proxy = default_tor;

    printf("URL: %s\n", url_arg);
    printf("User-Agent: %s\n", user_agent);
    (void)follow_redirects; (void)max_redirects;
    (void)head_request; (void)extra_headers;
    (void)post_data; (void)output_file; (void)auth_header;
\end{lstlisting}
\begin{lstlisting}[style=oldcode,language=C,firstnumber=last]
    return 0;
}
\end{lstlisting}
\FullSourceStop

The \texttt{(void)} casts silence unused-variable warnings until later steps
consume those fields. Parsing early means step~9 will only \emph{use} the
flags, not reinvent the CLI.
